% Part: methods
% Chapter: proofs
% Section: cant-do-it

\documentclass[../../../include/open-logic-section]{subfiles}

\begin{document}

\olfileid{mth}{prf}{cnt}

\olsection{મારાથી આ નહીં થાય!{}}

આપણને સૌને ક્યારેક હાર માની લેવાનું મન થાય છે. પણ તમે તે
\emph{કરી શકો છો}. તમારા અધ્યાપક અને શિક્ષણ-સહાયક તેમજ
સહાધ્યાયીઓ મદદ કરી શકે છે. તેમની મદદ માગો!{} મુશ્કેલી
ટાળવા અને હાર માનવાનું મન થાય ત્યારે શું કરવું તે માટે
આ કેટલીક સૂચનાઓ છે.

પ્રશ્નો સફળતાપૂર્વક ઉકેલી શકો તે માટે આ કરો:
\begin{enumerate}
\item \emph{શક્ય હોય એટલું વહેલું શરૂ કરો.} સત્ર દરમિયાન
  આપણે સૌ વ્યસ્ત થઈ જઈએ છીએ અને ઘણાને કામ ટાળવાની ટેવ
  સતાવે છે. તેથી ઘરકામ વહેલું શરૂ કરવું શ્રેષ્ઠ ઉપાયોમાંનો
  એક છે. એમ કરશો તો અટકો ત્યારે ઉકેલ શોધવા માટે સમય
  રહેશે (રડી પડવું ઉકેલ નથી).
\item \emph{સહાધ્યાયીઓ સાથે વાત કરો}. તમે એકલા નથી.
  વર્ગમાં બીજા લોકોને પણ મુશ્કેલી પડતી હશે---ભલે તેમની
  મુશ્કેલી જુદી હોય. તેમની સાથે વાત કરવાથી પ્રશ્નનો બીજો
  દૃષ્ટિકોણ મળે અને ઉકેલ સૂઝી શકે. અલબત્ત, તેમનો ઉકેલ
  સીધો ન ઉતારો: સંકેત માગો, અથવા તમે ક્યાં અટક્યા છો તે
  સમજાવીને આગળનું પગલું પૂછો. પછી તમને સમજ પડે ત્યારે
  તેમને પણ મદદ કરો. બીજાને સમજાવવાથી તમારી પોતાની
  સમજ પણ વધુ સ્પષ્ટ થાય છે.
\item \emph{મદદ માગો.} તમારા માટે ઘણા સ્રોતો છે---તમારા
  અધ્યાપક અને શિક્ષણ-સહાયક તમારી મદદ માટે છે અને તમે
  સફળ થાઓ એવું \emph{ઇચ્છે} છે. પ્રશ્ન ઉકેલવામાં અને
  કયા તબક્કે અટકો છો તે ઓળખવામાં તેઓ મદદ કરી શકશે.
\item \emph{થોડો વિરામ લો.} અટક્યા હો તો તેનું એક કારણ
  \emph{કદાચ} એ હોય કે તમે પ્રશ્નને બહુ લાંબો સમયથી જોઈ રહ્યા છો.
  નાનો વિરામ લો, ચા પીવો અથવા થોડો સમય બીજો પ્રશ્ન
  ઉકેલો; પછી તાજા મનથી પાછા આવો. એક રાત ઊંઘ્યા પછી
  ફરી વિચારો.
\end{enumerate}

ધ્યાન આપ્યું કે આ રીતો માટે સાબિતી પર ઘણું પહેલાંથી કામ
શરૂ કરેલું હોવું જોઈએ? સોંપવાની તારીખના આગલા દિવસે
વહેલી સવારે બે વાગ્યે શરૂ કર્યું હોય તો કદાચ તે બહુ
ઉપયોગી ન થાય.

આ નિરાશાજનક લાગશે, પણ સાબિતીનો પ્રશ્ન ઉકેલવો એવો પડકાર
છે જે અંતે સંતોષ આપે છે. કેટલાક લોકો આને વ્યવસાય તરીકે
કરે છે---એટલે તેમાં આનંદ લેવા જેવું કંઈક હશે જ. લગભગ
બીજી દરેક બાબતની જેમ પ્રશ્નો ઉકેલવા અને સાબિતીઓ કરવી
અભ્યાસ માગે છે. કેટલાક સહાધ્યાયીઓને આ સરળ લાગતું હશે;
તેમણે કદાચ પહેલેથી ઘણો અભ્યાસ કર્યો હશે. બહુ સહેલાઈથી
હાર ન માનો.

કોઈ પ્રશ્ન માટે સમય (અથવા ધીરજ) ખૂટી જાય તો તેમાં વાંધો
નથી. તેનો અર્થ તમે બુદ્ધિહીન છો કે તમને કદી નહીં સમજાય
એવો નથી. અધ્યાપક અથવા બીજા વિદ્યાર્થી પાસેથી ઉકેલની
રીત જાણો અને તમે ક્યાં ભૂલ્યા અથવા અટક્યા તે ઓળખો,
જેથી આગળ આવો પ્રશ્ન આવે ત્યારે એ ભૂલ ટાળી શકો. પછી
ઉકેલ જોયા વિના ફરી જાતે કરવાનો પ્રયાસ કરો. અને આગલી
વખતે વહેલું શરૂ કરો (અને વહેલું મદદ માગો).

\end{document}
